\appendix
\subsection{Video Generation Prompts}
\label{sec:vid_gen_prompts}

 For all real world tasks, the prompt to each video generation model consists of the initial RGB observation from one camera, as well as a language instruction of the form:

\begin{promptbox}[title={Real World Task Language Prompt}]

$<$TASK$>$ by one hand. The camera holds a still pose, not zooming in or out.

\end{promptbox}

\textbf{Values of \texttt{<TASK>}:}
\begin{itemize}
    \item \textit{Put Bread in Bowl:} The bread is grabbed and placed into the green bowl
    \item \textit{Open Oven:} The toaster oven is opened
    \item \textit{Cover Bowl:} The scarf is lifted by a corner and directly dragged over the blue bowl in one smooth motion without flinging the cloth or any dynamic / fast motions
    \item \textit{Pull Out Chair:} The chair is pulled out straight from under the table to the right by grabbing the middle
    \item \textit{Open Drawer:} The partially opened drawer is opened all the way out
    \item \textit{Sweep Pasta:} The brush with a green handle moves left to right to push the pasta into the compost bin
    \item \textit{Recycle Can:} The can is grabbed and dropped into the recycling bin
\end{itemize}

For the robustness evaluations, the word ``bread'' is replaced with ``donut'' or ``long piece of bread'' for the different object instances, but otherwise the prompt remains the same.

Since \algo leverages the position of a hand in the video to help select the grasp on the object, the ``by one hand'' phrase must be included. Additionally, ``the camera holds a still pose'' must be included to increase the probability that the generated videos do not have substantial camera motion, as the depth estimation pipeline assumes a still camera. 

\begin{promptbox}[title={Push-T Task Language Prompt}]

The T-shaped block slides and rotates smoothly to the center of the wooden platform.

\end{promptbox}

\begin{figure}[!thbp]
    \centering
    \includegraphics[width=0.96\linewidth]{figs/push_t_prompt.pdf}
    \caption{\textbf{Example Push-T image prompt.} To generate sufficiently accurate motions for the Push-T task, the visual component of the prompt includes both a start and end frame.}
    \label{fig:pusht_prompt}
    \vspace{-0.5em}
\end{figure}

For Push-T, since the task success requires accurate positions and to provide a better hint as to where the T block should go, we also provide a goal image of the T-block, as shown in Figure~\ref{fig:pusht_prompt}. At the time of evaluation, Veo 3 did not have the ability to take in an end frame, and hence was not considered for Push-T experiments. 

\begin{promptbox}[title={Open Door Task Language Prompt}]

The door opens all the way by itself to the right. The camera holds a still pose, not zooming in or out.

\end{promptbox}

The Open Door task prompt does not include a hand to perform the action as it is in a simulated environment. 

All prompts discussed in this section are identical for each video generation model evaluated. For Kling 2.1, a relevance value of 0.7 and a negative prompt of ``fast motion, morphing, camera motion'' are added.

\subsection{Particle Dynamics Model}
\label{sec:particle_dynamics_details}

\begin{figure}[!t]
    \centering
    \includegraphics[width=0.9\linewidth]{figs/particle_dynamics_model_fig.pdf}
    \caption{\textbf{Particle dynamics model.} The particle dynamics model primarily composed of a Point Transformer V3 backbone used in the Push-T task takes as input a set of feature-augmented particles and outputs delta position predictions for all particles in the scene.}
    \label{fig:particle_dynamics_model}
    \vspace{-0.5em}
\end{figure}

The particle dynamics model used in the Push-T task takes as input feature-augmented particles $\Tilde{x}_t \in \mathbb{R}^{N \times 14}$, consisting of the position, RGB value, normal vector, and push parameters and produces $\Delta \hat{x}_{t+1}$, the delta of positions of each point for the next timestep, as shown in Figure~\ref{fig:particle_dynamics_model}. The architecture of the particle dynamics model is a small Point Transformer V3~\cite{wu2024ptv3} backbone surrounded by MLPs to project between the desired input and output sizes. 

To get the input set of particles from the scene, there are 4 virtual cameras around the workspace, whose RGB-D observations are combined into a single point cloud. Then, points outside of the bounds of the wooden platform that the T-shaped block is on are discarded. Finally, this point cloud is voxel downsampled by randomly keeping only 1 particle per each 1.5cm sided cube. The same push parameter is appended to each particle's features. 

To train the particle dynamics model, we collect 500 transitions of random pushing actions. In this setting, we track particle positions before and after the push by using all objects' poses from the simulator for efficiency, but in practice this can also be tracked with CoTrackerV3~\cite{karaev2024cotracker3} and depth. 

\subsection{Push-T Planning Details}
\label{sec:pusht_planning}
To optimize the particle trajectory following cost with the learned particle based dynamics model, \algo replans after every single push with random shooting until the T-block is within the specified tolerance to the goal or a maximum number of pushes have been executed. In this setting, at each replanning step, $r$ push skill parameters are randomly sampled such that all pushes will make contact with the object of interest at different points and directions. Then, \algo selects the push skill parameter out of those $r$ ones that has the least cost according to the predicted particle positions from the dynamics model with respect to a subgoal of particle positions (may not necessarily be the final goal position). While the T-block is being pushed, \algo uses the online version of CoTrackerV3~\cite{karaev2024cotracker3} to track its motion, and once the push is completed, the tracked points are lifted into 3D, forming the updated particles of the T-block. We find that while parts of the T-block may be occluded during a pushing action, CoTrackerV3~\cite{karaev2024cotracker3} tends to recover visibility of previously occluded points when the gripper lifts up.

Since \algo tracks particles for the T-block while the particle dynamics model takes in downsampled particles of the entire scene, \algo performs nearest neighbor matching to find correspondences between the currently tracked particles and the feature-augmented particles which are the input to the dynamics model. With these correspondences, and the predicted delta positions of these corresponding particles, \algo computes the predicted particle positions of the T-block as $\hat{x}_{t+1} = \hat{x}_{t} + \Delta \hat{x}_{t+1}$.
 
For determining which timestep from the video should be used in the cost function as a subgoal, \algo first finds the timestep $t^{\star}$ where the tracked particles are closest to the particles in the \flow, as a single push can encompass the motion of many timesteps. \algo then chooses the particles from the timestep $\min(t^{\star} + L, t_{\text{end}})$ as the next subgoal for planning, where L is a fixed look ahead time amount, and $t_{\text{end}}$ is the final timestep of the video. In our experiments, we used $L=20$. We choose to use intermediate subgoals to be the target for random shooting as opposed to only the final particle positions of the T-block, as for cases involving substantial rotation, only having the final positions can result in the T-block have small translation error but large rotation error, eventually lead to timeouts.

\subsection{Grasp Selection}
\begin{figure}[!t]
    \centering
    \includegraphics[width=0.9\linewidth]{figs/grasp_selection.pdf}
    \caption{\textbf{Grasp selection with thumb position.} \algo selects a grasp closest to the position of the thumb in the video, if such a grasp exists within 2cm. The red sphere indicates the thumb position predicted by HaMer, with the red grasp being the selected one.}
    \label{fig:grasp_selection_thumb}
    \vspace{-0.5em}
\end{figure}

\algo uses AnyGrasp~\cite{fang2023anygrasp} to propose a set of up to 40 top-down grasps after applying the mask to the object of interest. Up to 20 grasps come from a point cloud in the coordinate frame of the camera, and another 20 grasps come from transforming those points into the coordinate frame of a virtual camera looking straight down in the center of the workspace. We added this addition virtual camera so that some more vertical grasps could be proposed. 

While for rigid objects consisting of one part, such as a piece of bread, it is typically fine to grasp at any stable position, for articulated objects, the part that moves needs to be grasped instead. To perform grasp selection in such cases, we exploit video generation models' ability to synthesize plausible hand-object interactions, typically where the hand first grabs the relevant part of the object and then proceeds to move it. \algo uses HaMer~\cite{pavlakos2024reconstructing} to detect the position of the hand, and in particular the thumb. If the predicted thumb position comes within 2cm of a proposed grasp, then such a grasp at the earliest timestep is chosen. An example of grasp selection with this method is shown in Figure~\ref{fig:grasp_selection_thumb}.

It is possible that there are no such grasps detected, either because the grasp planner proposes grasps that are not near the hand in the video or the detections from HaMer~\cite{pavlakos2024reconstructing} are not accurate enough. In such cases, \algo defaults to a heuristic of selecting the closest grasp to the points on the movable part of the rigid object (obtaining points on the movable part is described in the next section). 

\subsection{Movable Part Flow Filtering for Rigid-Grasp Dynamics Model}

\begin{figure*}[!ht]
    \centering
    \includegraphics[width=0.99\linewidth]{figs/in_the_wild_rollout.pdf}
    \caption{\textbf{In-the-Wild Task Rollouts.} The Franka successfully pulls out a chair, opens a partially opened drawer, sweeps pasta into the compost bin, and recycles an aluminum can.}
    \label{fig:in_the_wild_rollout}
    \vspace{-0.5em}
\end{figure*}


Since the rigid-grasp dynamics model assumes that points that are grasped move with the end-effector, it is necessary to identify what points are part of the movable object. To determine this, we employ the heuristic that individual flows that move at least 1 pixel on average per timestep are part of the movable flow. We empirically chose the 1 pixel threshold, as we noticed that points tracked on the stationary parts of articulated objects such as the oven had low averages of how many pixels they moved.

For the 2D tracks associated with the movable part, it may be possible that for certain intermediate frames they are 1 pixel off, making the individual tracks belong to the background or floor when lifted into 3D, causing part of the 3D object flow to be dramatically incorrect, leading to downstream execution failures (such as trying to push the toaster oven's door in the direction of the hinge). To remedy this issue, we utilize SAM 2~\cite{ravi2024sam2} with positive point prompts in the initial frame for the movable part and negative point prompts for the non-movable part, and use the part mask over each consecutive frame to constrain the 2D flow (if any tracked point is outside of the mask, it is considered to be invalid). 

\subsection{Real World Planning Details}
\label{sec:real_world_planning}

The optimization problem with a rigid-grasp assumption in the real world follows the same formulation as in Sec.~\ref{sec:problem}, where we optimize robot joint angles $\mathbf{q} \in \mathbb{R}^7$ to minimize:

\begin{equation}
\sum_{t=0}^{H-1} \lambda_{\text{task}}\big(\hat{x}^{\text{obj}}_t, \tilde{P}_t\big) + \lambda_{\text{control}}(\hat{x}_t, u_t)
\end{equation}

The control cost $\lambda_{\text{control}}(\hat{x}_t, u_t)$ is expanded into three components:
\begin{equation}
\lambda_{\text{control}}(\hat{x}_t, u_t) = w_r \mathcal{C}_r(\mathbf{q}_t) + w_s \mathcal{C}_s(\mathbf{q}_t, \mathbf{q}_{t-1}) + w_m \mathcal{C}_m(\mathbf{q}_t)
\end{equation}

where:
\begin{itemize}
\item $\mathcal{C}_r(\mathbf{q}_t)$: Reachability cost penalizing joint configurations outside the robot's workspace
\item $\mathcal{C}_s(\mathbf{q}_t, \mathbf{q}_{t-1})$: Pose smoothness cost measuring the difference between consecutive end-effector poses after running forward kinematics
\item $\mathcal{C}_m(\mathbf{q}_t)$: Manipulability cost encouraging configurations with good manipulability
\end{itemize}

The weight parameters are set to $w_r = 100$, $w_s = 1$, and $w_m = 0.01$ to encourage the robot to make smooth motions within its joint limits. The task cost $\lambda_{\text{task}}$ uses a weight of $w_f = 10$ and follows the same formulation as in Sec.~\ref{sec:problem}. The optimized end-effector poses after running forward kinematics are then fit with a B-spline and sampled such that each sampled pose is at least 1cm away from the previous and next pose and/or has a rotation difference of at least 20 degrees. To execute this planned trajectory, we use the IK solver from PyBullet~\cite{coumans2021pybullet} to get target joint angles and then a joint impedance controller from Deoxys~\cite{zhu2022viola} commands the Franka to reach those positions. 

\subsection{In-the-Wild Tasks}
\label{sec:in_the_wild_desc}

For In-the-Wild tasks, we consider:

\subsubsection{Pull Out Chair} 
In the Pull Out Chair task, a black rolling chair is placed underneath a table. 
A trial is considered a success if the chair is moved at least 5cm horizontally from its initial position. The reason for the limited motion requirement is that in certain configurations, it is difficult for the Franka robot to push the chair. 

\subsubsection{Open Drawer}
The Open Drawer task involves opening a partially opened drawer out to at least 90\% of its full possible extension. 

\subsubsection{Sweep Pasta}
In the Sweep Pasta task, there is a brush with a green handle placed against a wooden support structure, along with four pieces of dried pasta next to a compost bin. The objective is for the robot to grasp the handle of the brush and use the brush to push the pasta into the compost bin. If all pieces of pasta are inside of the compost bin, it is considered a success.  

\subsubsection{Recycle Can}
In the Recycle Can task, an aluminum can is placed in between a recycling and trash bin while upright. A trial is successful if the can is inside of the recycling bin.

See Fig.~\ref{fig:in_the_wild_rollout} for rollouts for each of these in-the-wild tasks.

\subsection{Open Door Reinforcement Learning Details}
\label{sec:open_door_rl}

For the Open Door task, we utilize Soft Actor-Critic (SAC)~\cite{haarnojaSAC2018} to train sensorimotor policies across three different embodiments: a Franka Panda, a Spot robot, and a GR1 humanoid arm. The policies are trained to follow the 3D object flow extracted from generated videos, which serves as a reward signal. 

\subsubsection{Hyperparameters}
The hyperparameters used for SAC training are detailed in Table~\ref{tab:sac_hyperparams}. These values were consistent across all reward types and embodiments with exception of the GR1, which uses 10000 training iterations due to the larger action space.

\begin{table}[h]
    \centering
    \begin{tabular}{lc}
        \toprule
        Hyperparameter & Value \\
        \midrule
        Learning rate & $3 \times 10^{-4}$ \\
        Discount factor ($\gamma$) & 0.99 \\
        Batch size & 256 \\
        Buffer size & $10^6$ \\
        Target update rate ($\tau$) & 0.005 \\
        Target network update freq & 1 \\
        Hidden layers & 2 \\
        Hidden units per layer & 256 \\
        Training iterations & 5000 \\
        Episode horizon & 500 \\
        \bottomrule
    \end{tabular}
    \caption{\textbf{SAC Hyperparameters for Open Door Task.}}
    \label{tab:sac_hyperparams}
\end{table}

\subsubsection{Reward Functions}
We compare policies trained with two different reward formulations: a handcrafted object state reward and a 3D object flow reward.

\textbf{Object State Reward:} The handcrafted reward $R_{\text{state}}$ consists of a reaching component $r_{\text{reach}}$ and a handle rotation component $r_{\text{rot}}$:
\begin{equation}
    r_{\text{reach}} = 0.25 \cdot (1 - \tanh(10 \cdot d_{\text{gripper, handle}}))
\end{equation}
\begin{equation}
    r_{\text{rot}} = \text{clip}\left(0.25 \cdot \frac{|\theta_{\text{handle}}|}{0.5\pi}, -0.25, 0.25\right)
\end{equation}
where $d_{\text{gripper, handle}}$ is the L2 distance between the robot's end-effector and the door handle, and $\theta_{\text{handle}}$ is the rotation angle of the handle. The total reward is $r_{\text{reach}} + r_{\text{rot}}$, unless the door hinge angle $\theta_{\text{hinge}} > 0.3$ radians, in which case a completion reward of 1.0 is returned.

\textbf{3D Object Flow Reward:} The 3D object flow reward $R_{\text{flow}}$ leverages the reference trajectory $P_{1:T}$ extracted from the video. It consists of a particle tracking term $r_{\text{particle}}$ and an end-effector alignment term $r_{\text{ee}}$:
\begin{equation}
    r_{\text{particle}} = 0.75 \cdot \frac{t^\star}{t_{\text{end}}}
\end{equation}
where $t^\star$ is the index of the closest timestep in the reference trajectory $P_{1:T}$ based on the current object particle positions $\hat{x}_t^{\text{obj}}$ in the door frame:
\begin{equation}
    t^\star = \argmin_{t \in \{1 \dots t_{\text{end}}\}} \frac{1}{n} \sum_{i=1}^n \|\hat{x}_t^{\text{obj}}[i] - P_t[i]\|_2
\end{equation}
The end-effector term $r_{\text{ee}}$ encourages the robot to stay near the object particles:
\begin{equation}
    r_{\text{ee}} = 0.25 \cdot (1 - \tanh(10 \cdot \|ee_{\text{door}} - \bar{x}\|_2))
\end{equation}
where $ee_{\text{door}}$ is the end-effector position in the door body frame and $\bar{x} = \frac{1}{n}\sum_{i=1}^n \hat{x}_t^{\text{obj}}[i]$ is the mean position of the object particles. The total reward is $R_{\text{flow}} = r_{\text{particle}} + r_{\text{ee}}$.

Instead of tracking the mean particle position with CoTrackerV3, we proceed to use a simpler approach of transforming the initial particles as the door angle changes for better computation efficiency. We thus consider the door joint angle to be a part of the state $x_t$.

\subsection{Video Generation Failures}
\begin{figure}[!thbp]
    \centering
    \includegraphics[width=0.96\linewidth]{figs/vid_gen_failures.pdf}
    \caption{\textbf{Video Generation Failure Examples.} (a) The bread undergoes object morphing, where it turns into a stack of crackers, causing the downstream tracking to fail. (b) Another green bowl appears due to model hallucination, causing a downstream execution failure where the bread is dropped onto the surface of the workspace instead of the original green bowl.}
    \label{fig:vid_gen_failures}
    \vspace{-0.5em}
\end{figure}

From the generated videos, we observe that there are two common failure modes: morphing and hallucination. Object morphing occurs when an existing object in the scene dramatically changes shape to something else, such as another object or an object with significantly different geometric properties than what it should be. Hallucination occurs when a new object (previously non-existent) appears in the scene. We show examples of morphing and hallucination for the Put Bread task in Fig.~\ref{fig:vid_gen_failures}.

\subsection{Limitations}
\algo has several limitations. First, it relies on a rigid-grasp assumption for real world manipulation, limiting the types of tasks that can be performed. While this work shows that a particle dynamics model can be used for other types of tasks such as non-prehensile pushing, training and scaling a particle dynamics model for the real world is non-trivial and can be considered for future work. Another limitation is that the total processing time to get 3D object flow depending on the video generation model is between 3 and 11 minutes, which limits its usability, with the main bottleneck being the video generation. Since \algo relies upon one angle in a generated video, it cannot handle heavy occlusions gracefully, such as when the human hand covers the majority of a small object. Future work may consider methods which can deal with such occlusions better, such as 3D point trackers~\cite{xiao2025spatialtracker} or full 4D representations.  
